En France, en 2021, une étude portant sur 100 personnes locataires du parc social (HLM) montre que :

	\begin{itemize}[label={\textbullet~}]
		\item 70 sont des personnes qui ne vivent pas en couple et parmi elles, la moitié est sans enfant ;
		\item Parmi les personnes vivant en couple, 12 n'ont pas d'enfant.
	\end{itemize}

	\begin{center}
		\begin{tblr}{|c|c|c|c|}
			\hline
			& sans enfant & au moins un enfant & Total \\
			\hline
			personne vivant en couple & 12 & 18 & 30 \\
			\hline
			personne ne vivant pas en couple & 35 & 35 & 70 \\
			\hline
			Total & 47 & 53 & 100 \\
			\hline
		\end{tblr}
	\end{center}

	On interroge une de ces personnes au hasard et on considère les évènements suivants :

	\begin{itemize}[label={\textbullet~}]
		\item $C$ : «\,la personne vit en couple\,»;
		\item $E$ : «\,la personne a au moins un enfant\,».
	\end{itemize}

	\begin{enumerate}
		\item Calculer la probabilité que la personne interrogée soit une personne sans enfant et ne vivant pas en couple.

		\item Calculer la probabilité de l'évènement $C \cap \overline{E}$ puis interpréter.

		\item Calculer la probabilité que la personne interrogée ne vive pas en couple sachant qu'elle a au moins un enfant.

		\item Dans cette question, on s'intéresse au nombre de pièces par logement dans le parc social.

		Une autre étude donne la répartition suivante :
		\begin{itemize}[label={\textbullet~}]
			\item $8\,\%$ des logements possèdent 5 pièces;
			\item $25\,\%$ ont 4 pièces;
			\item $40\,\%$ ont 3 pièces;
			\item $20\,\%$ ont deux pièces;
			\item et $7\,\%$ ont une seule pièce.
		\end{itemize}
		On choisit un logement au hasard et on note $X$ la variable aléatoire associant à chaque logement, son nombre de pièce(s).

		\begin{enumerate}
			\item Déterminer la loi de probabilité de $X$.

			\item Calculer $P(X \geqslant  2)$.

			Interpréter ce résultat dans le contexte de l'exercice.

			\item Montrer que $\mathrm{E}(X)=3,07$.

			Interpréter ce résultat.
		\end{enumerate}
	\end{enumerate}
